% Proofs for Section 13 (renormalization).
\section{Proofs for \texorpdfstring{Section~\ref{sec:ren}}{renormalization}}\label{app:ren}
% src: proofs of H4-L2, H4-T1, H4-P2 (odd-b step of HC_critic V8(i)), H4-X1, H4-T3, H4-X2, H4-P4, H4-T5, H4-P6, H4-P7, H4-T8, H4-P9, H4-X3, H4-P11 to H4-P17, H4-T18; H1-P1; H6-L1, H6-T2, H6-P11, H6-P12; written while drafting: prop:ren-exponents(b), the derivation for rem:ren-quartic, prop:ren-depth(d), thm:ren-closure(c) without equitability

The proofs follow the order of \cref{sec:ren}. On $\Z^D$ the symbol of $f$ is $\hat f(p)=\sum_xf(x)e^{-ip\cdot x}$, $p\in[-\pi,\pi]^D$, and Poisson summation for a subsampled function reads $\sum_{x'}f(bx')e^{-ip'\cdot x'}=b^{-D}\sum_{k\in\mathcal I_b}\hat f((p'+2\pi k)/b)$.

\begin{proof}[Proof of \cref{lem:ren-galerkin}]
(a) In the orthonormal basis $[\Phi,\Phi_\perp]$ the precision has the blocks $\Pi_\Phi$, $\Pi_{\perp\Phi}^\top$, $\Pi_{\perp\Phi}$ and $\Pi_{\perp\perp}$, and the block-inverse formula gives $(\Phi^\top\Sigma\Phi)^{-1}=\Pi_\Phi-\Pi_{\perp\Phi}^\top\Pi_{\perp\perp}^{-1}\Pi_{\perp\Phi}=\Pi_\Phi-\mathcal D$. As a congruence of $\Pi_{\perp\perp}^{-1}\succ0$, $\mathcal D\succeq0$, with $\mathcal D=0$ if and only if $\Pi_{\perp\Phi}=0$, and $\Pi_{\perp\perp}^{-1}\preceq\varkappa_\perp^{-1}I$ gives $\mathcal D\preceq\varkappa_\perp^{-1}\Pi_{\perp\Phi}^\top\Pi_{\perp\Phi}\preceq(\Vert\Pi_{\perp\Phi}\Vert^2/\varkappa_\perp)I$. The condition $\Pi_{\perp\Phi}=0$ says $\Pi\Phi=\Phi\Pi_\Phi$, that $\Pi$ maps $\operatorname{range}\Phi$ into itself. With uniform leak, $\Pi\one_Y=\varepsilon^\Pi\one_Y+L^\Pi\one_Y$, and $(L^\Pi\one_Y)_x$ is $-\sum_{y\in Y}a^\Pi_{xy}$ for $x\notin Y$ and $\sum_{z\notin Y}a^\Pi_{xz}$ for $x\in Y$; it is constant on every block for every $Y$ exactly when every member of $X$ has the same total coupling to $Y$, for all $X\neq Y$.

(b) $\mathcal D_{XY}=(\Pi_{\perp\Phi}\mathbf e_X)^\top\Pi_{\perp\perp}^{-1}(\Pi_{\perp\Phi}\mathbf e_Y)$, and $\Pi_{\perp\Phi}\mathbf e_X=\Phi_\perp^\top\Pi\phi_X$ has the norm $\lvert(I-\Phi\Phi^\top)\Pi\phi_X\rvert=\sigma_X$. The Cauchy--Schwarz inequality in the inner product of $\Pi_{\perp\perp}^{-1}\preceq\varkappa_\perp^{-1}I$ gives the bound, and $\sum_X\sigma_X^2=\Vert\Pi_{\perp\Phi}\Vert_F^2=\sigma^2_{\mathrm{cut}}$; the spectral norm is at least the norm of every column and at most the Frobenius norm. For uniform leak, $(I-\Phi\Phi^\top)$ annihilates $\varepsilon^\Pi\phi_X$ and is a contraction, so $\sigma_X\le\lvert L^\Pi\phi_X\rvert$. By the formula for $L^\Pi\one_X$ in (a), with $d_i$ the total coupling of agent $i$ across the boundary of $X$, to the other blocks for $i\in X$ and to $X$ for $i\notin X$, $\lvert L^\Pi\phi_X\rvert^2=\lvert X\rvert^{-1}\sum_id_i^2\le d_{\max}(X)\lvert X\rvert^{-1}\sum_id_i=2d_{\max}(X)\cut(X)/\lvert X\rvert$, since the members and the nonmembers each contribute $\cut(X)$ to $\sum_id_i$, and $\cut(X)/\lvert X\rvert=\xout(X)-\varepsilon^\Pi$.

(c) With $\mathsf n=(\lvert X\rvert^{1/2})_X$, $\Phi\mathsf n=\one$ and $\Pi\one=\varepsilon^\Pi\one$, so $\Pi_{\perp\Phi}\mathsf n=\varepsilon^\Pi\Phi_\perp^\top\one=0$ and $\mathcal D\mathsf n=0$. The block averages are $\bar e=\diag(\mathsf n)^{-1}y$, with exact precision $\diag(\mathsf n)(\Pi_\Phi-\mathcal D)\diag(\mathsf n)=\mathsf P\Pi\mathsf P^\top-\diag(\mathsf n)\mathcal D\diag(\mathsf n)$, and $\diag(\mathsf n)\mathcal D\diag(\mathsf n)\one=\diag(\mathsf n)\mathcal D\mathsf n=0$.

(d) For $X\neq Y$, $(\mathsf P\Pi\mathsf P^\top)_{XY}=\sum_{x\in X,y\in Y}\Pi_{xy}=-\cut(X,Y)\le0$, and $\sum_Y(\mathsf P\Pi\mathsf P^\top)_{XY}=\sum_{x\in X}(\Pi\one)_x=\sum_{x\in X}\varepsilon^\Pi_x\ge0$. The blocks are disjoint and nonempty, so $\mathsf P^\top u\neq0$ for $u\neq0$ and $u^\top\mathsf P\Pi\mathsf P^\top u>0$. For blocks of equal size $s$, $\Pi_\Phi=s^{-1}\mathsf P\Pi\mathsf P^\top$.
\end{proof}

\begin{proof}[Proof of \cref{thm:ren-decimation}]
The eliminated sites form the segments $\{kb+1,\dots,kb+b-1\}$, each coupled only to its two flanking retained sites, so $\Pi_{\bar{\mathcal R}\bar{\mathcal R}}$ is block diagonal with blocks $\mathsf S=a^\Pi\mathrm{tridiag}(-1,2\cosh m,-1)$ of size $b-1$, and the Schur complement couples only consecutive retained sites. The Dirichlet Green function is $(\mathsf S^{-1})_{ij}=\sinh(im)\sinh((b-j)m)/(a^\Pi\sinh m\sinh bm)$ for $1\le i\le j\le b-1$, extended symmetrically. As a function of $i$ it satisfies $2g_i\cosh m=g_{i-1}+g_{i+1}$ away from $i=j$ and vanishes at $i=0$ and $i=b$, and at $i=j$ the recursion leaves the remainder $[\sinh((j+1)m)\sinh((b-j)m)-\sinh(jm)\sinh((b-j-1)m)]/(\sinh m\sinh bm)=1$, by $\sinh A\sinh B=\frac12(\cosh(A+B)-\cosh(A-B))$. The coupling between the retained sites $0$ and $b$ is $(a^\Pi)^2(\mathsf S^{-1})_{1,b-1}=a^\Pi\sinh m/\sinh bm$, and each of the two adjacent segments lowers the diagonal entry by $(a^\Pi)^2(\mathsf S^{-1})_{11}=a^\Pi\sinh((b-1)m)/\sinh bm$, so $\Pi^{(1)}_{00}=2a^\Pi\cosh m-2a^\Pi\sinh((b-1)m)/\sinh bm$. Writing $\Pi^{(1)}_{00}=2a^{\Pi,(1)}\cosh m^{(1)}$ and using $\sinh((b-1)m)=\sinh bm\cosh m-\cosh bm\sinh m$ gives $\cosh m^{(1)}=\cosh bm$, and then $\varepsilon^{\Pi,(1)}=2a^{\Pi,(1)}(\cosh bm-1)=2a^\Pi\sinh m\tanh(bm/2)$. With $\cosh m=1+\mathfrak x/2$ this is \eqref{eq:ren-chebyshev}; $\mathsf{Ch}_b'(1)=b^2$ gives $\Lambda_{\mathrm{RG}}=b^2$ at $\mathfrak x=0$, and $m^{(l)}=b^lm$ makes $\mathfrak x^{(l)}=2(\cosh(b^lm)-1)$ increase to infinity for every $m>0$. The lengths: $\xi_{\mathrm{OZ}}=1/m$ and, for nearest-neighbor couplings, $\xi_2=1/(2\sinh(m/2))$ (\cref{prop:sca-lengths-relation}(c)); with $1/(4\sinh^2(y/2))=y^{-2}-1/12+O(y^2)$, $(\xi^{(1)}_2)^2=(bm)^{-2}-1/12+O(m^2)=\xi_2^2/b^2-(1-b^{-2})/12+O(m^2)$. Decimation is a marginalization, so the site variance is preserved, and indeed $1/(2a^{\Pi,(1)}\sinh m^{(1)})=1/(2a^\Pi\sinh m)$; finally $\sinh m/\sinh bm\to1/b$ as $m\to0$.
\end{proof}

\begin{proof}[Proof of \cref{prop:ren-ring}]
(a) Decimation evaluates $\Sigma$ at $bx'$, and $\mathsf p^{bx'}=(\mathsf p^b)^{x'}$.

(b) Write $\mathsf h=\sinh\mu$, $\mathsf g=\cosh\mu$, $\mathsf h'=\sinh\mu'$, $\mathsf g'=\cosh\mu'$ and $u=\cos p$. The symbols of $e^{-\mu\lvert x\rvert}$ and of $\mathsf p^xe^{-\mu'\lvert x\rvert}$ are $\mathsf h/(\mathsf g-u)$ and $\mathsf h'/(\mathsf g'-\mathsf pu)$, so
\[
\hat\Sigma=\frac{(\mathsf p\mathsf c_+\mathsf h+\mathsf c_-\mathsf h')(u_\ast-u)}{(\mathsf g-u)(\mathsf g'-\mathsf pu)},\qquad\hat\Pi=\frac{(\mathsf g-u)(\mathsf g'-\mathsf pu)}{(\mathsf p\mathsf c_+\mathsf h+\mathsf c_-\mathsf h')(u_\ast-u)}.
\]
Dividing the quadratic numerator by $u_\ast-u$ leaves an affine function of $u=\cos p$, a precision of range one, plus the constant $(\mathsf g-u_\ast)(\mathsf g'-\mathsf pu_\ast)$ times $1/(u_\ast-u)$, whose Fourier coefficients are proportional to $e^{-\arccosh(u_\ast)\lvert x\rvert}$ when $u_\ast>1$; the constant is nonzero because $u_\ast$ differs from $\mathsf g$ and from $\mathsf p\mathsf g'$, as follows. For $\mathsf p=1$, $u_\ast$ is a convex combination of $\mathsf g'$ and $\mathsf g$, which are distinct and above one. For $\mathsf p=-1$, which occurs for odd $b$, the function $x\mapsto\sinh(Bx)/\sinh x$ is increasing on $(0,\infty)$ for $B>1$, since its logarithmic derivative is $(\varphi(Bx)-\varphi(x))/x$ with $\varphi(y)=y\coth y$, and $\varphi'(y)=(\sinh y\cosh y-y)/\sinh^2y>0$, and $m'>m$ because $\cosh m'-\cosh m=c/a>0$. Hence, with $\mathsf c_-/\mathsf c_+=\sinh m/\sinh m'$,
\[
\frac{\mathsf c_-\mathsf h'}{\mathsf c_+\mathsf h}=\frac{\sinh(b^lm')/\sinh m'}{\sinh(b^lm)/\sinh m}>1,\qquad u_\ast-\mathsf g=\frac{\mathsf c_+\mathsf h(\mathsf g+\mathsf g')}{\mathsf c_-\mathsf h'-\mathsf c_+\mathsf h}>0,
\]
so $u_\ast>\mathsf g>1$ and the factor $\mathsf p\mathsf c_+\mathsf h+\mathsf c_-\mathsf h'$ is positive. The precision of passing is diffusive (\cref{thm:col-diffusive-precision}(e)), so its Kron reductions on rings whose length is a multiple of $b^l$ are diffusive (\cref{lem:col-kron}); on $\Z$ the entries are limits of those on rings, and nonpositive off-diagonal entries and nonnegative row sums survive the limit.

(c) By (a) the staggered term after $l$ steps is $\mathsf c_-\mathsf p^xe^{-b^lm'\lvert x\rvert}$; for $\lvert x\rvert\ge1$ its ratio to $\mathsf c_+e^{-b^lm\lvert x\rvert}$ is at most $(\mathsf c_-/\mathsf c_+)e^{-b^l(m'-m)}$, and at $x=0$ it is $\mathsf c_-$. A pure exponential $\mathsf c_+e^{-\mu\lvert x\rvert}$ is the covariance of a nearest-neighbor precision with mass $\mu$ (\cref{thm:ren-decimation}), and $\sinh m/\sinh m'\to0$ as $\varepsilon\to0$, since $m\to0$ while $m'$ stays positive.
\end{proof}

\begin{proof}[Proof of the first part of \cref{ex:ren-two-dim}]
Odd sites have only even neighbors, so $\Pi_{\mathrm{oo}}=(\varepsilon^\Pi+4a^\Pi)I$ and the Schur complement onto the even sites is $(\varepsilon^\Pi+4a^\Pi)I-(a^\Pi)^2N_{\mathrm{eo}}N_{\mathrm{oe}}/(\varepsilon^\Pi+4a^\Pi)$, where $(N_{\mathrm{eo}}N_{\mathrm{oe}})_{xy}$ counts the odd sites adjacent to both $x$ and $y$: four for $x=y$, two for $y-x\in\{(\pm1,\pm1)\}$, one for $y-x\in\{(\pm2,0),(0,\pm2)\}$, and none otherwise. In the rotated coordinates the first displacements are nearest neighbors and the second diagonal neighbors, and the leak is $(\varepsilon^\Pi+4a^\Pi)-16(a^\Pi)^2/(\varepsilon^\Pi+4a^\Pi)=\varepsilon^\Pi(\varepsilon^\Pi+8a^\Pi)/(\varepsilon^\Pi+4a^\Pi)$.
\end{proof}

\begin{proof}[Proof of \cref{thm:ren-flows}]
(a) Decimation is Poisson summation. For averaging, $\bar e_{x'}=b^{-D}\sum_{v\in\mathcal I_b}e_{bx'+v}$, so $\Cov(\bar e_0,\bar e_{x'})=g(bx')$ with $g(x)=b^{-2D}\sum_{u,v\in\mathcal I_b}\Sigma(x+v-u)$, the convolution of $\Sigma$ with the autocorrelation of the box, whose symbol is $b^{-2D}\lvert\sum_{v\in\mathcal I_b}e^{ip\cdot v}\rvert^2=F_b(p)$; subsampling $g$ gives the formula.

(b) Decimation. By Poisson summation in the transverse coordinates, the sum of the decimated covariance over the coarse hyperplane $\{x'_1=t'\}$ is $b^{-(D-1)}\sum_{k_\perp}\tilde\Sigma(bt',2\pi k_\perp/b)$, with $k_\perp\in\{0,\dots,b-1\}^{D-1}$ and $\tilde\Sigma(t,p_\perp)=\sum_{x_\perp}\Sigma(t,x_\perp)e^{-ip_\perp\cdot x_\perp}$, the covariance in $t$ of the one-dimensional field with symbol $\hat\Pi(\cdot,p_\perp)$. For the nearest-neighbor precision, $\hat\Pi(p_1,p_\perp)=\varepsilon^\Pi+2a^\Pi\sum_{d\ge2}(1-\cos p_d)+2a^\Pi(1-\cos p_1)$, so $\tilde\Sigma(t,p_\perp)$ is a pure exponential in $\lvert t\rvert$ with rate $m(p_\perp)$, $\cosh m(p_\perp)=1+(\varepsilon^\Pi+2a^\Pi\sum_{d\ge2}(1-\cos p_d))/(2a^\Pi)$, which exceeds $m_1=m(0)$ for every $p_\perp=2\pi k_\perp/b$ with $k_\perp\neq0$. For passing, $\hat\Pi=(1-\hat W^2)/q$ with $\hat W=c+2a\sum_d\cos p_d$, and $\tilde\Sigma(\cdot,p_\perp)$ is the covariance of \cref{thm:sca-oz}(e) with $c$ replaced by $c(p_\perp)=c+2a\sum_{d\ge2}\cos p_d$: its smooth rate solves $\cosh m=(1-c(p_\perp))/(2a)$, which is smallest at $p_\perp=0$, where it is $m_1$, and its staggered rate solves $\cosh m'(p_\perp)=(1+c(p_\perp))/(2a)$, above $\cosh m_1=(1-c-2a(D-1))/(2a)$ for every $p_\perp$ because $c>0$. In both cases the term $k_\perp=0$, a positive multiple of $e^{-m_1\lvert t\rvert}$ plus a faster staggered part, decays at the rate $m_1$ per site, hence $bm_1$ per coarse step, and every other term decays faster. Averaging. Summed over the coarse transverse positions and the transverse offsets, every fine transverse displacement occurs exactly $b^{D-1}$ times, so the hyperplane sum of the averaged field is $b^{-D-1}\sum_{u_1,v_1=0}^{b-1}\tilde\Sigma(bt'+v_1-u_1,0)$. For $t'\ge1$ every argument is positive, so the sum is a positive multiple of $e^{-bm_1t'}$ plus a term of the faster rate $bm'_1$. The effective ratio follows from $\cosh m_1=1+\varepsilon/(2a)$, with $\cosh(bm_1)-1=b^2(\cosh m_1-1)(1+O(m_1^2))$.

(c) Expand at $p'=0$. For decimation,
\[
\hat\Sigma^{\mathsf{Dec}}(0)=b^{-D}(\hat\Sigma(0)+\mathsf S_1),\qquad-\partial^2_{p'_d}\hat\Sigma^{\mathsf{Dec}}(0)=b^{-D-2}(-\partial^2_{p_d}\hat\Sigma(0)+\mathsf S_{2,d}),
\]
with $-\partial^2_{p_d}\hat\Sigma(0)=2\xi^2_{2,d}\hat\Sigma(0)$ by \cref{prop:sca-lengths-relation}(a); the covariance-side length is half the ratio of the two. The decimated precision is diffusive, being a Kron reduction on tori (\cref{lem:col-kron}) and a limit of those on $\Z^D$. Its symbol $1/\hat\Sigma^{\mathsf{Dec}}$ is twice differentiable, because $\hat\Pi$ is (finite second moments) and $\hat\Pi\ge\varepsilon^\Pi>0$, so by Fatou's lemma its nonnegative couplings have finite second moments, and \cref{prop:sca-lengths-relation}(a) identifies its second-moment length with the covariance-side one. Averaging: $F_b(2\pi k/b)=0$ for $k\neq0$, so $\hat\Sigma^{\mathsf{Avg}}(0)=b^{-D}\hat\Sigma(0)$. With $\sin^2(p/2)/(b^2\sin^2(p/(2b)))=1-p^2(1-b^{-2})/12+O(p^4)$, the term $k=0$ contributes $\hat\Sigma(0)[1-p'^2_d((1-b^{-2})/12+\xi^2_{2,d}/b^2)]$ in coordinate $d$ up to fourth order; a term $k=k_d\mathbf e_d$ contributes $(p'^2_d/4)\hat\Sigma(2\pi k_d\mathbf e_d/b)/(b^2\sin^2(\pi k_d/b))+O(p'^4)$; and terms with a nonzero coordinate other than $d$ carry the factor $\sin^2(p'_{d'}/2)$ and do not enter $\partial^2_{p'_d}$ at $0$. Collecting terms with $\hat\Sigma(0)=1/\varepsilon^\Pi$ gives the formula. For the constants, decimation gives $(\xi^{\mathsf{Dec}}_{2,d})^2-\xi^2_{2,d}/b^2=(\varepsilon^\Pi\mathsf S_{2,d}-\varsigma_d^2\mathsf S_1)/(2b^2(1+\varepsilon^\Pi\mathsf S_1))$, using $2\xi^2_{2,d}\varepsilon^\Pi=\varsigma_d^2$, which tends to $-\varsigma_d^2\mathsf S^{(0)}_1/(2b^2)$ because $\mathsf S_1$ and $\mathsf S_{2,d}$ stay bounded as $\varepsilon^\Pi\to0$, $\hat\Pi-\varepsilon^\Pi$ being positive at the points $2\pi k/b$, $k\neq0$; for averaging the last term is $O(\varepsilon^\Pi)$ for the same reason. For the nearest-neighbor kernel in one dimension, $\mathsf S^{(0)}_1=\sum_{k=1}^{b-1}(4a^\Pi\sin^2(\pi k/b))^{-1}=(b^2-1)/(12a^\Pi)$ and $\varsigma^2=2a^\Pi$, which give $-(1-b^{-2})/12$, and the classical sum $\sum_{k=1}^{b-1}\sin^{-4}(\pi k/b)=(b^4+10b^2-11)/45$ shows that the last term of the averaging formula is then $(\varepsilon^\Pi/a^\Pi)(b^4+10b^2-11)/(720b^2)+O((\varepsilon^\Pi)^2)$.
\end{proof}

\begin{proof}[Proof of \cref{ex:ren-block-average}]
For $X\ge1$ every argument $bX+v-u$ with $u,v\in\{0,\dots,b-1\}$ is positive, so $\Sigma^{\mathsf{Avg}}(X)=b^{-2}\mathsf c_\Sigma e^{-bmX}\sum_{u,v}e^{-m(v-u)}$, and $\sum_{u,v}e^{-m(v-u)}=\sum_ue^{mu}\sum_ve^{-mv}=\mathsf g_b$. At $X=0$ the sum is $\sum_{u,v}e^{-m\lvert v-u\rvert}=b+2\sum_{j=1}^{b-1}(b-j)e^{-jm}$, which is positive and smaller than $\mathsf g_b$, because pairs with $u>v$ contribute $e^{-m(u-v)}<e^{m(u-v)}$; this gives $\mathsf k_b\in(0,1)$. With $\sum_Xe^{-bm\lvert X\rvert}e^{-ipX}=\sinh bm/(\cosh bm-\cos p)$, the symbol of $\Sigma^{\mathsf{Avg}}$ is a positive multiple of $(\mathsf e_0+\mathsf k_b\cos p)/(\cosh bm-\cos p)$. It is positive for every $p$, so $\mathsf e_0>\mathsf k_b$, and the precision symbol is a positive multiple of $(\cosh bm-\cos p)/(\mathsf e_0+\mathsf k_b\cos p)$. The Fourier coefficients of $1/(\mathsf e_0+\mathsf k_b\cos p)$ are $(-\mathsf h_b)^{\lvert n\rvert}/(\mathsf e_0^2-\mathsf k_b^2)^{1/2}$, where $\mathsf h_b\in(0,1)$ solves $(1+\mathsf h_b^2)/(2\mathsf h_b)=\mathsf e_0/\mathsf k_b$, by the geometric series. Multiplying by $\cosh bm-\cos p$ gives, for $\lvert n\rvert\ge1$, the coefficient $(-\mathsf h_b)^{\lvert n\rvert}[\cosh bm+(\mathsf h_b^{-1}+\mathsf h_b)/2]/(\mathsf e_0^2-\mathsf k_b^2)^{1/2}=-(-\mathsf h_b)^{\lvert n\rvert-1}[\mathsf h_b\cosh bm+(1+\mathsf h_b^2)/2]/(\mathsf e_0^2-\mathsf k_b^2)^{1/2}$, which is $-\mathsf c_1(-\mathsf h_b)^{\lvert n\rvert-1}$ with $\mathsf c_1>0$ once the positive multiple is included. The covariance symbol is a positive ratio of trigonometric polynomials of degree one, $(\mathsf e_0+\mathsf k_b\cos p)\propto\lvert1+\mathsf h_be^{ip}\rvert^2$ over $\cosh bm-\cos p\propto\lvert1-e^{-bm}e^{ip}\rvert^2$, so the block averages form an autoregressive moving-average field of order $(1,1)$. As $m\to0$, $\mathsf g_b\to b^2$ and $\mathsf g_b-\sum_{u,v}e^{-m\lvert v-u\rvert}=\sum_{u>v}2\sinh(m(u-v))=mb(b^2-1)/3+O(m^3)$, so $\mathsf k_b=m(b^2-1)/(3b)+O(m^2)$, $\mathsf e_0=bm-\mathsf k_b+O(m^2)$ and $\mathsf k_b/\mathsf e_0\to\mathsf t_b$, which gives the limit of $\mathsf h_b$; $\mathsf t_2=1/3$ gives $3-2\sqrt2=0.1716$, and $\mathsf t_b\to1/2$ gives $2-\sqrt3$.
\end{proof}

\begin{proof}[Proof of \cref{prop:ren-exponents}]
(a) By \cref{thm:ren-flows}(c), $(\xi^{(1)}_{2,d})^2=\xi^2_{2,d}/b^2+c_{b,D}+O(\varepsilon^\Pi)$ for both maps, with $c_{b,D}$ independent of the leak and $\varepsilon^\Pi=\varsigma_d^2\mathfrak x_{2,d}/2$; inverting gives $\mathfrak x^{(1)}_{2,d}=b^2\mathfrak x_{2,d}/(1+O(b^2\mathfrak x_{2,d}))$. If every leak of a finite closed population vanished, $\Pi\one=0$ and $\Pi$ would be singular, contradicting \textup{(M)}.

(b) Near $q=0$, $\hat\Sigma(q)=1/\hat\Pi(q)=1/\varepsilon^\Pi-\mathsf a_\sigma\lvert q\rvert^\sigma/(\varepsilon^\Pi)^2+O(\lvert q\rvert^{\min(\sigma',2\sigma)})$. Away from $2\pi\Z^D$, $\hat\Sigma$ is twice continuously differentiable, so $\mathsf S(p')=\sum_{k\in\mathcal I_b\setminus\{0\}}\hat\Sigma((p'+2\pi k)/b)$ is twice continuously differentiable near $p'=0$; it is even, because $\hat\Sigma$ is even and $2\pi$-periodic and $k\mapsto-k$ permutes the nonzero classes modulo $b$, so $\mathsf S(p')=\mathsf S_1+O(\lvert p'\rvert^2)$. Hence $b^D\hat\Sigma^{\mathsf{Dec}}(p')=1/\varepsilon^\Pi+\mathsf S_1-b^{-\sigma}\mathsf a_\sigma\lvert p'\rvert^\sigma/(\varepsilon^\Pi)^2+O(\lvert p'\rvert^{\min(\sigma',2\sigma,2)})$, and inverting,
\[
\hat\Pi^{\mathsf{Dec}}(p')=\frac{b^D\varepsilon^\Pi}{1+\varepsilon^\Pi\mathsf S_1}+\frac{b^{D-\sigma}\mathsf a_\sigma}{(1+\varepsilon^\Pi\mathsf S_1)^2}\lvert p'\rvert^\sigma+O\bigl(\lvert p'\rvert^{\min(\sigma',2\sigma,2)}\bigr).
\]
For averaging, each term $k\neq0$ carries the factor $F_b((p'+2\pi k)/b)=O(\lvert p'\rvert^2)$, because some coordinate has $k_d\neq0$ and $\sin^2((p'_d+2\pi k_d)/2)=\sin^2(p'_d/2)$, and the term $k=0$ carries $F_b(p'/b)=1+O(\lvert p'\rvert^2)$; so $b^D\hat\Sigma^{\mathsf{Avg}}(p')=1/\varepsilon^\Pi-b^{-\sigma}\mathsf a_\sigma\lvert p'\rvert^\sigma/(\varepsilon^\Pi)^2+O(\lvert p'\rvert^{\min(\sigma',2\sigma,2)})$, which inverts to the stated coefficients. The ratio $\varepsilon^\Pi/\mathsf a_\sigma$ is therefore multiplied by $b^\sigma(1+\varepsilon^\Pi\mathsf S_1)$ or by $b^\sigma$, and $\mathsf S_1$ stays bounded as $\varepsilon^\Pi\to0$ because $\hat\Pi-\varepsilon^\Pi$ is positive at the points $2\pi k/b$, $k\neq0$; so $\Lambda_{\mathrm{RG}}=b^\sigma$ at the massless point. If the couplings had a finite second moment, then $\hat\Pi(p)-\varepsilon^\Pi=\sum_xa^\Pi(x)(1-\cos p\cdot x)\le\frac12\lvert p\rvert^2\sum_xa^\Pi(x)\lvert x\rvert^2$, contradicting $\mathsf a_\sigma\lvert p\rvert^\sigma\gg\lvert p\rvert^2$ as $p\to0$; so $\varsigma^2=\infty$ and $\xi_2=\infty$.

(c) Under the reversible diffusion the stationary autocovariance at lag $u$ has the symbol $e^{-u\hat\Pi(p)}\hat\Sigma(p)$, and decimation acts on it by the alias sum of \eqref{eq:ren-symbols}, so the term $k$ relaxes at $\hat\Pi((p'+2\pi k)/b)$. For $k=0$, $\hat\Pi(p'/b)=\varepsilon^\Pi+\sum_xa^\Pi(x)(1-\cos(p'\cdot x/b))=\varepsilon^\Pi+\sum_xa^\Pi(x)(p'\cdot x)^2/(2b^2)+o(\lvert p'\rvert^2b^{-2})$, by dominated convergence, since $\lvert1-\cos\theta-\theta^2/2\rvert\le\theta^2\min\{1,\theta^2/24\}$ and the couplings have finite second moments; the quadratic term is $\sum_d\varsigma_d^2p'^2_d/(2b^2)$ when the mixed moments vanish. In the time unit $b^2$ the rate is $b^2\varepsilon^\Pi+\sum_xa^\Pi(x)(p'\cdot x)^2/2+o(\lvert p'\rvert^2)$; under (b) the same holds with $b^\sigma$ and $\mathsf a_\sigma\lvert p'\rvert^\sigma$. For $k\neq0$ and $p'\in[-\pi,\pi]^D$, a coordinate with $k_d\neq0$ has $(p'_d+2\pi k_d)/b\in[\pi/b,2\pi-\pi/b]$, so for the nearest-neighbor kernel $\hat\Pi-\varepsilon^\Pi\ge2a^\Pi(1-\cos(\pi/b))=4a^\Pi\sin^2(\pi/(2b))$.
\end{proof}

\begin{proof}[Derivation for \cref{rem:ren-quartic}]
At $\varepsilon^\Pi=0$ write $u=\hat\Pi(p'/b)=\alpha_2p'^2/b^2+\alpha_4p'^4/b^4+O(p'^6)$ and $\mathsf S(p')=\mathsf S^{(0)}_1+O(p'^2)$ as in the proof of \cref{prop:ren-exponents}(b). Then $\hat\Pi^{\mathsf{Dec}}=b^Du/(1+u\mathsf S)=b^D(u-u^2\mathsf S^{(0)}_1+O(p'^6))$, whose coefficients of $p'^2$ and $p'^4$ are $b^{-2}\alpha_2$ and $b^{-4}(\alpha_4-\alpha_2^2\mathsf S^{(0)}_1)$, up to the factor $b^D$. For the nearest-neighbor kernel $\alpha_2=a^\Pi$, $\alpha_4=-a^\Pi/12$ and $\mathsf S^{(0)}_1=(b^2-1)/(12a^\Pi)$, so the ratio is $b^{-2}(-1/12-(b^2-1)/12)=-1/12$.
\end{proof}

\begin{proof}[Proof of \cref{thm:ren-nested}]
(a) For $x$ in a level-1 block $X$ and another level-1 block $Y$, $\sum_{y\in Y}a^{\mathrm H}(x,y)=ba^{\mathrm H}_{l(X,Y)}$ does not depend on $x$, with $l(X,Y)$ the level of the smallest common block. Hence $\Pi\one_Y$ is constant on every level-1 block, $\operatorname{range}\Phi$ is $\Pi$-invariant and $\Pi_{\perp\Phi}=0$, so $\mathcal D=0$ (\cref{lem:ren-galerkin}), and $\mathrm dy=-\Phi^\top\Pi(\Phi y+\Phi_\perp\Phi_\perp^\top x)\mathrm dt+\sqrt2\Phi^\top\mathrm d\mathsf B_t$ with $\Phi^\top\Pi\Phi_\perp=\Pi_{\perp\Phi}^\top=0$. Symmetric weights with the nested structure have block sums constant on level-1 blocks, so the partition is lumpable for them (\cref{def:met-lumpable}), and \cref{thm:ren-closure}(b) gives a closed recursion for the block averages.

(b) $(\Pi_\Phi)_{XY}=b^{-1}\sum_{x\in X,y\in Y}\Pi_{xy}=-ba^{\mathrm H}_l$ for blocks whose smallest common block has level $l$, which is level $l-1$ in the coarse medium, and $ba^{\mathrm H}_l=(\mathfrak qA_{\mathrm h})\mathfrak q^{(l-1)-1}/((b-1)b^{(l-1)-1})$; the row sums are $b^{-1}\sum_{x\in X}(\Pi\one)_x=\varepsilon^\Pi$. So the amplitude becomes $\mathfrak qA_{\mathrm h}$ and $\mathfrak x_l$ is multiplied by $1/\mathfrak q$. If lengths grow by $b^{1/D_e}$ per level, $1/\mathfrak q=(b^{1/D_e})^{y_{\mathrm{RG}}}$ defines $y_{\mathrm{RG}}$, and $D_e-y_{\mathrm{RG}}=D_e\ln(b\mathfrak q)/\ln b=\alpha_{\mathrm H}$ by \cref{prop:abs-hierarchy}(c); the massless lattice field has the Green function $\lvert x\rvert^{2-D_e}$.

(c) Each member of a level-$(l+1)$ block $C$ has total coupling $A_{\mathrm h}\mathfrak q^{l'-1}$ to its level-$l'$ cousins, so $\cut(C)/\lvert C\rvert=A_{\mathrm h}\sum_{l'=l+2}^{n_{\mathrm H}}\mathfrak q^{l'-1}$. Inside $C$, a vector that is constant on the level-$(j-1)$ sub-blocks of one level-$j$ block, sums to zero over it and vanishes elsewhere in $C$ is an eigenvector of $L_C$ with eigenvalue $g_j=a^{\mathrm H}_jb^j+\sum_{l'=j+1}^{l+1}a^{\mathrm H}_{l'}b^{l'-1}(b-1)=A_{\mathrm h}\mathfrak q^{j-1}b/(b-1)+\sum_{l'=j+1}^{l+1}A_{\mathrm h}\mathfrak q^{l'-1}$: members of the same sub-block contribute nothing, those of other sub-blocks of the same level-$j$ block contribute $a^{\mathrm H}_jb^j$ because the vector sums to zero there, and those outside the level-$j$ block contribute their couplings. These vectors, over $j=1,\dots,l+1$, span $\one_C^\perp$, and $g_j-g_{j+1}=A_{\mathrm h}\mathfrak q^{j-1}(b-\mathfrak q)/(b-1)>0$, so $\lambda_2(L_C)=g_{l+1}=bA_{\mathrm h}\mathfrak q^l/(b-1)$, and with uniform leak $r(C)=(\varepsilon^\Pi+\lambda_2(L_C))/(\varepsilon^\Pi+\cut(C)/\lvert C\rvert)$. By (b) the level-$k$ medium is nested with amplitude $A_{\mathrm h}\mathfrak q^k$, in which $C$ is a block of level $l+1-k$, and the same formula gives the same value; for equal blocks the orthonormal and the block-sum conventions give the same contrast (\cref{lem:ren-galerkin}(d)).

(d) As $n_{\mathrm H}\to\infty$ the denominator of \eqref{eq:ren-nested-contrast} is $\varepsilon^\Pi+A_{\mathrm h}\mathfrak q^{l+1}/(1-\mathfrak q)$, and dividing by $A_{\mathrm h}\mathfrak q^l$ gives the formula. After clearing denominators $r_l\ge r_0$ reads $\mathfrak x_l(r_0-1)\le b/(b-1)-r_0\mathfrak q/(1-\mathfrak q)$, whose right side is positive exactly when $r_0<r_\infty$; since $\mathfrak x_l=\mathfrak x_0\mathfrak q^{-l}$ increases with $l$, the solutions are $l\le l^\ast$.
\end{proof}

\begin{proof}[Proof of \cref{prop:ren-depth}]
(a) With uniform leak \cref{thm:col-contrast-size} gives $\Var_d(S)\le\xi^2_{2,d}/(r_0-1)$ in every coordinate, so $\sum_d\Var_d(S)\le D\max_d\xi^2_{2,d}/(r_0-1)$, while $\sum_d\Var_d(S)=(2\lvert S\rvert^2)^{-1}\sum_{i,j\in S}\lvert x_i-x_j\rvert^2\ge(\lvert S\rvert-1)/(2\lvert S\rvert)\ge\frac14$, since distinct lattice points are at distance at least one.

(b) By \cref{thm:ren-decimation} the decimated precision is nearest-neighbor with mass $b^lm$, so $\xi^{(l)}_2=1/(2\sinh(b^lm/2))\le1/(b^lm)$, because $\sinh y\ge y$. The condition of (a), $1/(4\sinh^2(b^lm/2))\ge(r_0-1)/4$, is $b^l\le(2/m)\arcsinh((r_0-1)^{-1/2})$, and $\arcsinh y\le y$.

(c) In one dimension $\xi^{(l)}_2=1/(2\sinh(b^lm/2))\le1/(2b^l\sinh(m/2))=\xi_2/b^l$ by (b), because $\sinh$ is convex on $[0,\infty)$ and vanishes at $0$. Insert $\xi^{(l)}_{2,d}\le\xi_{2,d}/b^l$ into (a). For passing, $\xi_2\le(a/\varepsilon)^{1/2}$ (\cref{prop:sca-lengths-relation}(c)).

(d) Let $X_0=\{0,\dots,B-1\}^D$, and for $x\in X_0$ let $k_d(x,z)=\lfloor(x_d+z_d)/B\rfloor$ be the block displacement in coordinate $d$ of the pair $(x,x+z)$. By \cref{lem:ren-galerkin}(d) the coarse medium is translation invariant on the coarse lattice, with leak $B^D\varepsilon^\Pi$ and second moment $\sum_{x\in X_0}\sum_za^\Pi(z)k_d(x,z)^2$ in coordinate $d$. For $z_d\ge0$ write $z_d=\mathsf jB+s$ with $0\le s<B$; then $k_d=\mathsf j+\one\{x_d\ge B-s\}$, so $\sum_{x_d=0}^{B-1}k_d^2=(B-s)\mathsf j^2+s(\mathsf j+1)^2=(z_d^2+s(B-s))/B$, and the same holds for $z_d<0$ by symmetry. The other $D-1$ coordinates of $x$ contribute the factor $B^{D-1}$, so the second moment is $B^{D-2}\sum_za^\Pi(z)(z_d^2+s_d(z)(B-s_d(z)))$, and dividing by twice the leak gives the formula. The inequality is $s(B-s)\le(B-1)z_d^2$: for $1\le\lvert z_d\rvert<B$, $s=\lvert z_d\rvert$ and $B-\lvert z_d\rvert\le(B-1)\lvert z_d\rvert$, with equality only at $\lvert z_d\rvert=1$; for $\lvert z_d\rvert\ge B$, $s(B-s)\le B^2/4<(B-1)z_d^2$; for $z_d=0$ both sides vanish. Block indicators of level $l+1$ are sums of those of level $l$, so coarse media compose and the level-$l$ medium is that of cubic blocks of side $b^l$; then $(\xi^{(l)}_{2,d})^2\le\xi^2_{2,d}/b^l$, with equality for nearest-neighbor couplings, and (a) gives $b^l\le4D\max_d\xi^2_{2,d}/(r_0-1)$.
\end{proof}

\begin{proof}[Proof of \cref{prop:ren-exchangeable}]
In the exchangeable medium a region of $n'$ agents has $\lambda_2=n'\bar a$ and $\cut/n'=(N-n')\bar a$, so its contrast $(\varepsilon^\Pi+n'\bar a)/(\varepsilon^\Pi+(N-n')\bar a)$ increases with $n'$, and the supremum is the whole population, $1+N\bar a/\varepsilon^\Pi$. With uniform leak $r(C)=(\varepsilon^\Pi+\lambda_2(L_C))/(\varepsilon^\Pi+\mathsf c)$, $\mathsf c=\cut(C)/n$, and since $\varepsilon^{\max}_C/\bar\varepsilon_C=1<\bar r_{\mathrm{hom}}$, $C$ is a collective at some admissible level exactly when $r(C)>\bar r_{\mathrm{hom}}$, that is $\varepsilon^\Pi(\varepsilon^\Pi+\lambda_2(L_C))>(\varepsilon^\Pi+\mathsf c)(\varepsilon^\Pi+N\bar a)$, or $\varepsilon^\Pi(\lambda_2(L_C)-\mathsf c-N\bar a)>\mathsf cN\bar a\ge0$, which is \eqref{eq:ren-exchangeable}. The vector $\mathbf e_i-\one_C/n$ is orthogonal to $\one_C$, has squared norm $(n-1)/n$ and quadratic form $d_C(i)$, so $\lambda_2(L_C)\le\frac n{n-1}\min_id_C(i)$, and $\lambda_2(L_C)>N\bar a=\frac N{N-1}\bar d$ gives the last claim.
\end{proof}

\begin{proof}[Proof of \cref{prop:ren-per-level}]
By \cref{thm:ren-nested}(b) and \cref{lem:ren-galerkin}(d), the level-$l$ coarse medium is, up to a positive factor, the nested medium on $M=b^m$ units, $m=n_{\mathrm H}-l$, with amplitude $A_l=A_{\mathrm h}\mathfrak q^l$ and leak $\varepsilon^\Pi$. Its level-1 blocks have $\lambda_2=bA_l/(b-1)$ and $\mathsf c=\cut/b=A_l(S_m-1)$ (\cref{thm:ren-nested}(c)), every unit has total coupling $A_lS_m$, and the exchangeable comparison has the pair coupling $A_lS_m/(M-1)$. By \cref{prop:ren-exchangeable} the blocks are collectives if and only if $\varepsilon^\Pi(\lambda_2-\mathsf c-MA_lS_m/(M-1))>\mathsf cMA_lS_m/(M-1)$, and $\lambda_2-\mathsf c-MA_lS_m/(M-1)=A_l\Delta_{\mathrm H}(b,\mathfrak q,m)$ and $\mathsf cMA_lS_m/(M-1)=A_l^2MS_m(S_m-1)/(M-1)$; this is (b). For $m=1$ the medium is exchangeable, hence its own comparison medium, and no region exceeds its baseline.

(a) The level-1 blocks are (b) at $l=0$. Let $C$ be a block of level $l+1\ge2$, judged at level 0. If $l+1\le n_{\mathrm H}-1$, then $\lambda_2(L_C)\le g_{l+1}=bA_{\mathrm h}\mathfrak q^l/(b-1)$, an eigenvalue of $L_C$ on $\one_C^\perp$ by the proof of \cref{thm:ren-nested}(c), while $\cut(C)/\lvert C\rvert\ge A_{\mathrm h}\mathfrak q^{l+1}$ and $N\bar a>\bar d=A_{\mathrm h}S_{n_{\mathrm H}}\ge A_{\mathrm h}\sum_{j=0}^{l+1}\mathfrak q^j$, so $\lambda_2(L_C)-\cut(C)/\lvert C\rvert-N\bar a<A_{\mathrm h}[\mathfrak q^l/(b-1)-2\mathfrak q^{l+1}-1]\le A_{\mathrm h}[\mathfrak q^l(1-2\mathfrak q)-1]<0$ for every $\mathfrak q>0$. If $C$ is the whole population, $\lambda_2(L)\le\frac N{N-1}\bar d=N\bar a$ by the degree bound of \cref{prop:ren-exchangeable}. In both cases \eqref{eq:ren-exchangeable} fails.

(c) $\Delta_{\mathrm H}=\frac{2b-1}{b-1}-S_m\bigl(2+\frac1{M-1}\bigr)$. For $m=2$, $\Delta_{\mathrm H}>0$ reads $\frac b{b-1}-\frac{b^2}{b^2-1}>\mathfrak q\bigl(1+\frac{b^2}{b^2-1}\bigr)$, that is $\mathfrak q<b/(2b^2-1)$. For $m\ge3$, $S_m<1/(1-\mathfrak q)$ and $1/(M-1)\le1/(b^3-1)$, so $\Delta_{\mathrm H}>\frac{2b-1}{b-1}-\frac{2b^3-1}{(1-\mathfrak q)(b^3-1)}$. At $\mathfrak q=b/(2b^2-1)$, $1-\mathfrak q=(2b+1)(b-1)/(2b^2-1)$, and the right side is nonnegative if and only if $(2b^2-1)(2b^3-1)\le(4b^2-1)(b^3-1)$, which is $b^2(b-2)\ge0$; it increases as $\mathfrak q$ decreases. So $\mathfrak q<b/(2b^2-1)$ gives $\Delta_{\mathrm H}>0$ for every $m\ge2$, and $m=2$ shows the converse. The second statement combines this with (b).

(d) For $m\ge2$, $S_m\ge1+\mathfrak q$ and $2+1/(M-1)>2$, so $\Delta_{\mathrm H}<\frac{2b-1}{b-1}-2-2\mathfrak q\le\frac{2b-1}{b-1}-2-\frac2b=\frac{2-b}{b(b-1)}\le0$ when $\mathfrak q\ge1/b$; then (a) and (b) exclude every level.

(e) $r_l-1=(\lambda_2-\mathsf c)/(\varepsilon^\Pi+\mathsf c)\le\lambda_2/\varepsilon^\Pi=bA_{\mathrm h}\mathfrak q^l/((b-1)\varepsilon^\Pi)$, while the baseline of the level-$k$ medium satisfies $\bar r^{(k)}_{\mathrm{hom}}-1=\frac{M_k}{M_k-1}A_{\mathrm h}\mathfrak q^kS_{n_{\mathrm H}-k}/\varepsilon^\Pi>A_{\mathrm h}\mathfrak q^k/\varepsilon^\Pi$, and $b\mathfrak q^{l-k}/(b-1)\le b\mathfrak q/(b-1)<1/(b-1)\le1$ for $l>k$ when $\mathfrak q<1/b$. A level admissible at level $k$ exceeds $\bar r^{(k)}_{\mathrm{hom}}$, hence every $r_l$ with $l>k$.

(f) Put $\mathfrak q_\ast=1/(2b-1)$. For $m\ge2$, $S_m$ increases strictly with $\mathfrak q$, so $\Delta_{\mathrm H}=\frac{2b-1}{b-1}-2S_m-\frac{S_m}{M-1}$ decreases strictly with $\mathfrak q$, and for $\mathfrak q<1$, $2S_m=2(1-\mathfrak q^m)/(1-\mathfrak q)$ gives $\Delta_{\mathrm H}=\frac{2b-1}{b-1}-\frac2{1-\mathfrak q}+\frac{2\mathfrak q^m}{1-\mathfrak q}-\frac{S_m}{M-1}$, which tends to $\frac{2b-1}{b-1}-\frac2{1-\mathfrak q}$ as $m\to\infty$, a limit that is positive exactly when $\mathfrak q<\mathfrak q_\ast$. At $\mathfrak q_\ast$, $2/(1-\mathfrak q_\ast)=(2b-1)/(b-1)$, so $\Delta_{\mathrm H}(b,\mathfrak q_\ast,m)=\frac{(2b-1)^{1-m}}{b-1}-\frac{S_m}{b^m-1}$; since $2b-1\ge b+1$, $(2b-1)^{m-1}\ge(b+1)^{m-1}\ge\sum_{j=0}^{m-1}b^j=(b^m-1)/(b-1)$, and $S_m>1$, this is negative for every $m\ge2$, and so is $\Delta_{\mathrm H}(b,\mathfrak q,m)$ for every $\mathfrak q\ge\mathfrak q_\ast$; then (a) and (b) exclude every level. For the increase in $m$, let $0<\mathfrak q\le\mathfrak q_\ast$, $m\ge2$ and $M=b^m$. Then
\[
\Delta_{\mathrm H}(b,\mathfrak q,m+1)-\Delta_{\mathrm H}(b,\mathfrak q,m)=\frac{S_m}{M-1}-\frac{S_{m+1}}{bM-1}-2\mathfrak q^m=\frac{(b-1)MS_m-\mathfrak q^m(M-1)(2bM-1)}{(M-1)(bM-1)},
\]
and $S_m/\mathfrak q^m=\sum_{j=1}^m\mathfrak q^{-j}\ge(2b-1)^m+(2b-1)^{m-1}=2b(2b-1)^{m-1}$, so the numerator is at least $\mathfrak q^m[2b(b-1)(2b-1)^{m-1}M-(M-1)(2bM-1)]$. For $m=2$ the bracket equals $(b-1)[2b^3(b-2)+b+1]>0$. For $m\ge3$, $(b-1)(2b-1)^{m-1}/b^m=\frac{b-1}{2b-1}\bigl(\frac{2b-1}b\bigr)^m\ge\frac{(b-1)(2b-1)^2}{b^3}=4-\frac8b+\frac5{b^2}-\frac1{b^3}\ge\frac98$, the last expression being increasing in $b\ge2$, so the bracket is at least $\frac94bM^2-(M-1)(2bM-1)>0$. Hence $\Delta_{\mathrm H}$ increases strictly with $m\ge2$ when $\mathfrak q\le\mathfrak q_\ast$ (which gives the negativity at $\mathfrak q_\ast$ a second time, the limit there being zero). For $\mathfrak q<\mathfrak q_\ast$ the $m\ge2$ with $\Delta_{\mathrm H}(b,\mathfrak q,m)>0$ are therefore exactly those with $m\ge m^\ast$, with $m^\ast=2$ exactly when $\mathfrak q<b/(2b^2-1)$ by (c), and (b) with $m=n_{\mathrm H}-l$ gives the statement.
\end{proof}

\begin{proof}[Proof of \cref{thm:ren-memory}]
(a) is \cref{lem:ren-galerkin}(a). (b) is \cref{thm:agy-mori-zwanzig}(b) in the coordinates $(y,\Phi_\perp^\top x)$, in which the drift has the blocks $\Pi_\Phi$, $\Pi_{\perp\Phi}^\top$, $\Pi_{\perp\Phi}$ and $\Pi_{\perp\perp}$; $\Gamma(u)$ is a congruence of $e^{-\Pi_{\perp\perp}u}\succeq0$, $\Vert e^{-\Pi_{\perp\perp}u}\Vert=e^{-\varkappa_\perp u}$, and $\int_0^\infty e^{-\Pi_{\perp\perp}u}\mathrm du=\Pi_{\perp\perp}^{-1}$.

(c) By the Courant--Fischer theorem, $\varpi_M\le\max$ of the Rayleigh quotient over $\operatorname{range}\Phi$, which is $\varkappa_\Phi$, and $\varpi_{M+1}\ge\min$ over $\operatorname{range}\Phi_\perp$, which is $\varkappa_\perp$. From $\Pi\mathsf V_{\mathrm f}=\mathsf V_{\mathrm f}\diag(\varpi_{\mathrm f})$ and $\Pi\Phi=\Phi\Pi_\Phi+\Phi_\perp\Pi_{\perp\Phi}$, the matrix $\mathsf Q_\Theta=\mathsf V_{\mathrm f}^\top\Phi$ solves $\diag(\varpi_{\mathrm f})\mathsf Q_\Theta-\mathsf Q_\Theta\Pi_\Phi=\mathsf V_{\mathrm f}^\top\Phi_\perp\Pi_{\perp\Phi}$. With the midpoint $\mathsf c_{\mathrm m}=(\varkappa_\perp+\varkappa_\Phi)/2$, $\mathsf M_1=\diag(\varpi_{\mathrm f})-\mathsf c_{\mathrm m}I$ and $\mathsf M_2=\mathsf c_{\mathrm m}I-\Pi_\Phi$ are at least $\frac12(\varkappa_\perp-\varkappa_\Phi)$, the equation reads $\mathsf M_1\mathsf Q_\Theta+\mathsf Q_\Theta\mathsf M_2=\mathsf V_{\mathrm f}^\top\Phi_\perp\Pi_{\perp\Phi}$, and its solution $\mathsf Q_\Theta=\int_0^\infty e^{-t\mathsf M_1}\mathsf V_{\mathrm f}^\top\Phi_\perp\Pi_{\perp\Phi}e^{-t\mathsf M_2}\mathrm dt$ has norm at most $\Vert\Pi_{\perp\Phi}\Vert/(\varkappa_\perp-\varkappa_\Phi)$; $\Vert\mathsf Q_\Theta\Vert=\sin\Theta$. The coordinates $z=\mathsf V^\top x$ obey $\mathrm dz=-\diag(\varpi)z\mathrm dt+\sqrt2\mathrm d(\mathsf V^\top\mathsf B)$ and are independent, and $y=\mathsf Tz_{\mathrm s}+\Phi^\top\mathsf V_{\mathrm f}z_{\mathrm f}$; the smallest singular value of $\mathsf T$ is $\cos\Theta>0$. $\Cov(y_{\mathrm f})=\mathsf Q_\Theta^\top\diag(\varpi_{\mathrm f})^{-1}\mathsf Q_\Theta\preceq\varkappa_\perp^{-1}\mathsf Q_\Theta^\top\mathsf Q_\Theta$, and $\Cov(y)^{-1}=\Pi_\Phi-\mathcal D\preceq\varkappa_\Phi I$, which gives the relative weight. For the drift, $\Phi^\top\Pi\mathsf V_{\mathrm s}=\mathsf T\diag(\varpi_{\mathrm s})$ and $\Phi^\top\Pi=\Pi_\Phi\Phi^\top+\Pi_{\perp\Phi}^\top\Phi_\perp^\top$ give $\mathsf T\diag(\varpi_{\mathrm s})\mathsf T^{-1}-\Pi_\Phi=\Pi_{\perp\Phi}^\top(\Phi_\perp^\top\mathsf V_{\mathrm s})\mathsf T^{-1}$, with $\Vert\Phi_\perp^\top\mathsf V_{\mathrm s}\Vert=\sin\Theta$ and $\Vert\mathsf T^{-1}\Vert=1/\cos\Theta$.

(d) For $f$ orthogonal to the block constants, $f^\top\Pi f\ge\varepsilon^\Pi\lvert f\rvert^2+\sum_jf_{C_j}^\top L_{C_j}f_{C_j}\ge\xin^{\min}\lvert f\rvert^2$, since the couplings across blocks add a positive semidefinite Laplacian and each $f_{C_j}$ is orthogonal to $\one_{C_j}$, hence zero on a region with one member, so that only the regions with at least two members enter the minimum. For $f=\Phi\mathsf y$, $f^\top L^\Pi f=\sum_{j<k}\cut(C_j,C_k)(\mathsf y_j/n_j^{1/2}-\mathsf y_k/n_k^{1/2})^2\le2\sum_j(\cut(C_j)/n_j)\mathsf y_j^2$, and $\cut(C_j)/n_j=\xout(C_j)-\varepsilon^\Pi$. The spectral norm is at most the Frobenius norm. For common rates, $\mathsf t_\Theta\le\Vert\Pi_{\perp\Phi}\Vert/(\xin-2\xout)$ and $\varkappa_\Phi/\varkappa_\perp\le2\xout/\xin=2/r$.
\end{proof}

\begin{proof}[Proof of \cref{prop:ren-perron}]
$\Pi_{C_jC_j}$ has nonpositive off-diagonal entries and a connected graph, so $\mathsf c_{\mathrm p}I-\Pi_{C_jC_j}$ is nonnegative and irreducible for large $\mathsf c_{\mathrm p}$, and its Perron vector, the eigenvector of the smallest eigenvalue of $\Pi_{C_jC_j}$, is positive.

(a) $\Pi_{\perp\Phi^{\mathsf P}}=0$ if and only if $\Pi\Phi^{\mathsf P}\mathbf e_j\in\operatorname{range}\Phi^{\mathsf P}$ for every $j$. The $C_j$ block of $\Pi\Phi^{\mathsf P}\mathbf e_j$ is $\Pi_{C_jC_j}\mathsf u_j$, proportional to $\mathsf u_j$, and its $C_k$ block is $\Pi_{C_kC_j}\mathsf u_j$. If it vanishes, $\mathcal D=0$ and $\Gamma=0$ by \cref{thm:ren-memory}(a),(b), and $\mathfrak f$ reduces to $\sqrt2\Phi^{\mathsf P\top}\mathrm d\mathsf B_t$, the increment of a standard Brownian motion, so AR0 is exact. Conversely AR0 has stationary covariance $\Pi_{\Phi^{\mathsf P}}^{-1}$, which equals $\Cov(y)=(\Pi_{\Phi^{\mathsf P}}-\mathcal D^{\mathsf P})^{-1}$ only if $\mathcal D^{\mathsf P}=0$, that is $\Pi_{\perp\Phi^{\mathsf P}}=0$.

(b) The first statements are \cref{thm:ren-memory} in the frame $\Phi^{\mathsf P}$, whose complement inside each particle holds that particle's other modes. A stationary process $\mathrm dy=-\mathsf A_{\mathrm M}y\mathrm dt+\sqrt2\mathsf G_{\mathrm M}^{1/2}\mathrm d\mathsf W_t$ has autocovariance $e^{-\mathsf A_{\mathrm M}u}\Cov(y)$, while the exact autocovariance $\Phi^{\mathsf P\top}e^{-\Pi u}\Pi^{-1}\Phi^{\mathsf P}$ has derivative $-\Phi^{\mathsf P\top}\Phi^{\mathsf P}=-I$ at $u=0$. Matching requires $\mathsf A_{\mathrm M}\Cov(y)=I$, and the Lyapunov equation $\mathsf A_{\mathrm M}\Cov(y)+\Cov(y)\mathsf A_{\mathrm M}^\top=2\mathsf G_{\mathrm M}$ then gives $\mathsf G_{\mathrm M}=I$.

(c) $\Pi\one=\varepsilon^\Pi\one$ and $\Pi\succeq\varepsilon^\Pi I$, so $\one/N^{1/2}$ is a slowest mode; it is a combination of block indicators and lies in $\operatorname{range}\Phi^{\mathsf P}$ only if every $\mathsf u_j$ is proportional to $\one_{C_j}$. The Rayleigh quotient gives $\lambda_{\min}(\Pi_{\Phi^{\mathsf P}})\ge\varepsilon^\Pi$. Let $\lambda=\varepsilon^\Pi$, $\Pi v=\lambda v$ and $v=\Phi^{\mathsf P}\mathsf a+\Phi^{\mathsf P}_\perp\mathsf b$. Since $\lambda\le\lambda_{\min}(\Pi_{\Phi^{\mathsf P}})<\varkappa_\perp$, the second block row gives $\mathsf b=-(\Pi_{\perp\perp}-\lambda)^{-1}\Pi_{\perp\Phi^{\mathsf P}}\mathsf a$ with $\mathsf a\neq0$, and the first gives $[\Pi_{\Phi^{\mathsf P}}-\Pi_{\perp\Phi^{\mathsf P}}^\top(\Pi_{\perp\perp}-\lambda)^{-1}\Pi_{\perp\Phi^{\mathsf P}}]\mathsf a=\lambda\mathsf a$; hence $\lambda\ge\lambda_{\min}(\Pi_{\Phi^{\mathsf P}})-\Vert\Pi_{\perp\Phi^{\mathsf P}}\Vert^2/(\varkappa_\perp-\lambda)$, and $\varkappa_\perp-\lambda\ge\varkappa_\perp-\lambda_{\min}(\Pi_{\Phi^{\mathsf P}})$. For block averages, $\Pi_\Phi\mathsf n=\Phi^\top\Pi\one=\varepsilon^\Pi\mathsf n$ with $\mathsf n=(\lvert X\rvert^{1/2})_X$.

(d) $(\Pi_\Phi)_{XX}=\xout(X)=\lvert X\rvert^{-1}\sum_{x\in X}(\varepsilon^\Pi_x+a^\partial_x)$, so for $M$ blocks of $N/M$ agents $M^{-1}\tr\Pi_\Phi=N^{-1}\sum_i(\varepsilon^\Pi_i+a^\partial_i)\le N^{-1}\sum_i\Pi_{ii}$, with equality only if no agent is coupled inside its block; $\mathcal D\succeq0$ gives the second inequality. A block average has variance $\phi_X^\top\Sigma\phi_X/\lvert X\rvert\le\lambda_{\max}(\Sigma)/\lvert X\rvert$.
\end{proof}

\begin{proof}[Proof of \cref{ex:ren-no-contrast}]
The coarse medium of blocks of $b$ consecutive sites has the coupling $\cut(X,X\pm1)=a^\Pi$, one edge, and the leak $b\varepsilon^\Pi$ (\cref{lem:ren-galerkin}(d)), so its second moment in block units is $2a^\Pi$ and its squared length $a^\Pi/(b\varepsilon^\Pi)$. For the nearest-neighbor kernel $\xi^2_2=a^\Pi/\varepsilon^\Pi$ (\cref{sec:col-size}), and \cref{thm:ren-flows}(c) gives the exact length $\xi^2_2/b^2+O(1)$. For the recursion, $W^{(1)}=\diag(\lvert X\rvert)^{-1}\mathsf PW\mathsf P^\top$ has self weight $c+2a(b-1)/b$ and neighbor weight $a/b$. A nearest-neighbor coarse network with message weight $\mathsf a_{\mathrm c}$ and leak $\varepsilon$ per round has the axis rate $\mathsf m_{\mathrm c}$ with $\cosh\mathsf m_{\mathrm c}=1+\varepsilon/(2\mathsf a_{\mathrm c})$ (\cref{thm:sca-oz}(c)); equating it with $bm$ gives $\mathsf a_{\mathrm c}=\varepsilon/(2(\cosh bm-1))$, and $\cosh bm-1=b^2(\cosh m-1)(1+O(m^2))=b^2\varepsilon(1+O(\varepsilon))/(2a)$.
\end{proof}

\begin{proof}[Proof of \cref{prop:ren-adiabatic}]
(a) Substitute $x_{\mathrm{int}}=\mathsf C_{\mathrm{ir}}x_{\mathrm{ret}}+x_\perp$ into $\mathrm dx_{\mathrm{ret}}=-(M_{\mathrm{rr}}x_{\mathrm{ret}}+M_{\mathrm{ri}}x_{\mathrm{int}})\mathrm dt+\sqrt2\mathsf N_{\mathrm r}\mathrm d\mathsf B_{\mathrm r}$. At stationarity $\Cov(x_\perp,x_{\mathrm{ret}})=\Sigma_{\mathrm{ir}}-\mathsf C_{\mathrm{ir}}\Sigma_{\mathrm{rr}}=0$.

(b) Itô's rule for the linear map $x_\perp=x_{\mathrm{int}}-\mathsf C_{\mathrm{ir}}x_{\mathrm{ret}}$ gives $\mathrm dx_\perp=\mathrm dx_{\mathrm{int}}-\mathsf C_{\mathrm{ir}}\mathrm dx_{\mathrm{ret}}$; substituting both rows and $x_{\mathrm{int}}=\mathsf C_{\mathrm{ir}}x_{\mathrm{ret}}+x_\perp$ and collecting terms gives the stated drift, the lower block row of $\mathcal TM\mathcal T^{-1}$ with $\mathcal T=\bigl(\begin{smallmatrix}I&0\\-\mathsf C_{\mathrm{ir}}&I\end{smallmatrix}\bigr)$. By variation of constants $x_\perp(t)=e^{-M_{\perp\perp}t}x_\perp(0)-\int_0^te^{-M_{\perp\perp}(t-t')}M_{\perp\mathrm r}x_{\mathrm{ret}}(t')\mathrm dt'+\sqrt2\int_0^te^{-M_{\perp\perp}(t-t')}(\mathsf N_{\mathrm i}\mathrm d\mathsf B_{\mathrm i}-\mathsf C_{\mathrm{ir}}\mathsf N_{\mathrm r}\mathrm d\mathsf B_{\mathrm r})(t')$, and inserting it into (a) gives the integro-differential equation with $\Gamma_{\mathrm F}$. By (a), $x_\perp$ is the Mori complement of \cref{thm:agy-mori-zwanzig}(a).

(c) If $M_{\mathrm{ri}}=0$ the retained rows do not contain $x_{\mathrm{int}}$, so $x_{\mathrm{ret}}$ solves an autonomous linear equation and is Markov. If $M_{\mathrm{ri}}z\neq0$ for some $z$, start from $x_{\mathrm{ret}}(0)=0$ with $x_{\mathrm{int}}(0)=0$ or $x_{\mathrm{int}}(0)=z$: the mean of $x_{\mathrm{ret}}(t)$ is $[e^{-Mt}]_{\mathrm{ri}}z=-tM_{\mathrm{ri}}z+O(t^2)$ in the second case and zero in the first, so no transition law of $x_{\mathrm{ret}}$ alone reproduces both. Active states are retained, so $M_{\mathrm{ai}}\neq0$ gives $M_{\mathrm{ri}}\neq0$.
\end{proof}

\begin{proof}[Proof of \cref{prop:ren-codes}]
Substitute $m_j=\eta+\bar e_J+\tilde e_j$ and use $\sum_j\mathsf w_jH_j=H^{\mathrm{agg}}_J$ and $\sum_j\mathsf w_j\tilde e_j=0$. (a) $\rank T_c\bar U_J\le\rank T_c=r_c$, and the stacked code has row space $\sum_j\operatorname{row}(H_j)$. (b) With a common code $H$, $\sum_j\mathsf w_jH\tilde e_j=0$. With uniform weights the Cauchy--Schwarz inequality gives $\lvert\sum_jn^{-1}(H_j-H^{\mathrm{agg}}_J)\tilde e_j\rvert^2\le n^{-1}\sum_j\Vert H_j-H^{\mathrm{agg}}_J\Vert^2\cdot n^{-1}\sum_j\lvert\tilde e_j\rvert^2$. The fluctuations of $\tilde e_j$ about their world-conditional means are those of the connected parts, and $\sum_j\Var(\tilde e_j\mid\eta)=K\tr((I-\one\one^\top/n)\Sigma^{\mathrm{sc}}_{JJ})\le K(n-1)/\xin(J)$, because every unit vector orthogonal to $\one_J$ has variance at most $1/\xin(J)$ (\cref{prop:col-coherence-average}(a)). (c) The term $H\bar e_J$ has world-conditional variance $H^2\Var(\phi_J^\top e)/n\ge H^2/(n\xout(J))$ by \cref{prop:col-coherence-average}(a), the channel noises are independent of it with variance $R_c/n$, and the second moment is at least the variance. The level-0 receiver assigns the variance $R_c/n$, so its believed precision $n/R_c$ exceeds the true precision $n/(H^2/\xout(J)+R_c)$ at least by the stated factor.
\end{proof}

\begin{proof}[Proof of \cref{prop:ren-code-levels}]
The strong law of large numbers and the variance of an average of independent matrices give the first statements, and singular values are $1$-Lipschitz in the operator norm (Weyl's inequality). (a) Left invariance of the Haar law gives $V\E U=\E U$ for every rotation $V$, and for $K\ge2$ the only vector fixed by every rotation is $0$; apply this to each column. With $\E U=0$, $\E\Vert T_c\bar U^{(l)}\Vert_F^2=b^{-l}\E\Vert T_cU\Vert_F^2=b^{-l}\Vert T_c\Vert_F^2$, since the Frobenius norm is invariant under right multiplication by an orthogonal matrix. (b), (c) $\E\sin\vartheta=0$ by symmetry, which gives $\E U$; the code with rows $\mathbf e_1^\top$ and $\mathbf e_2^\top$ then has $T_c\E U$ with singular values $(1,\lvert\E\cos\vartheta\rvert)$ for rotations about $\mathbf e_1$ and $(\lvert\E\cos\vartheta\rvert,\lvert\E\cos\vartheta\rvert)$ for rotations about $\mathbf e_3$. (d) A message translated with the true translation has the identity as its effective frame (\cref{sec:fw}).
\end{proof}

\begin{proof}[Proof of \cref{prop:ren-relay-rank}]
(a) $x^\top\bar Px=\sum_j\mathsf w_j\Vert P_{V_j}x\Vert^2\le\Vert x\Vert^2$, with equality if and only if $x\in V_j$ for every $j$, and for $0\preceq\bar P\preceq I$ the equality $\bar Px=x$ is equivalent to $x^\top\bar Px=\Vert x\Vert^2$. Also $x^\top\bar Px=0$ if and only if every $P_{V_j}x=0$, so $\ker\bar P=(\sum_jV_j)^\perp$. (b) $\bar P^n$ is symmetric positive semidefinite with eigenvalues $\lambda_k(\bar P)^n$, which are its singular values, and the ratios $(\lambda_k/\lambda_1)^n$ do not increase with $n$. $V_\cap$ is the eigenspace of the eigenvalue one, $V_\cap^\perp$ is invariant, and on it the largest eigenvalue $\lambda_\ast$ is below one by (a), so $\bar P^n-P_\cap=(\bar P(I-P_\cap))^n$ has norm $\lambda_\ast^n$. The effective rank tends to the multiplicity of $\lambda_1$, which is $\dim V_\cap$ when $V_\cap\neq\{0\}$; one meta-level hop has rank $\dim\sum_jV_j$.
\end{proof}

\begin{proof}[Proof of \cref{prop:ren-hayek}]
Each aggregate $\psi^{(l+1)}_C$ lies in the span of its summands, so the span of the functionals of level $l+1$ lies in that of level $l$, and its dimension is at most the number of its generators. The sufficiency statements are \cref{prop:abs-hayek} applied to each decision, with a public message whose rows span the functionals of level $l+1$, respectively those of the members.
\end{proof}

\begin{proof}[Proof of \cref{prop:ren-literal-levels}]
(a) $I+\mathrm{Adj}$ of a connected $\Delta$-regular graph is primitive, with Perron root $1+\Delta$ and Perron vector $\one/n^{1/2}$, so $(I+\mathrm{Adj})^{n_{\mathrm c}}/(1+\Delta)^{n_{\mathrm c}}\to\one\one^\top/n$ as $n_{\mathrm c}\to\infty$. After $n_{\mathrm c,1}$ rounds every member of group $J$, the representative included, holds natural parameters $(1+\Delta_1)^{n_{\mathrm c,1}}n^{-1}(\one^\top h_J(0),\one^\top\Lambda_J(0))(1+o(1))$, with the same factor for every group. At level two \cref{prop:agy-literal-consensus}(a) on a regular graph gives the limit mean $\sum_J\one^\top h_J(0)/\sum_J\one^\top\Lambda_J(0)$, in which every member contributes $\Lambda_\circ m_\circ$ and $\Lambda_\circ$ once. Induction on the number of levels gives the product, and the believed precisions grow geometrically at every level by \cref{prop:agy-literal-consensus}(b). (b) For a connected group, $(I+\mathrm{Adj}_J)^{n_{\mathrm c,1}}=\rho(I+\mathrm{Adj}_J)^{n_{\mathrm c,1}}\mathsf v^{(J)}\mathsf v^{(J)\top}(1+o(1))$, and the level-two limit mean is the ratio of the entering natural parameters weighted by the level-two Perron vector (\cref{prop:agy-literal-consensus}(a)). (c) Contributions add exactly, so the top level holds $(\Lambda_\circ m_\circ+\sum_ih^s_i,\Lambda_\circ+\sum_i\Lambda^s_i)$, the natural parameters of the union posterior, as in \cref{prop:agy-consensus}.
\end{proof}

\begin{proof}[Proof of \cref{prop:ren-lrg}]
$e^{-\tau_{\mathrm L}(\varepsilon^\Pi I+L^\Pi)}=e^{-\tau_{\mathrm L}\varepsilon^\Pi}e^{-\tau_{\mathrm L}L^\Pi}$, and the scalar factor cancels in the normalization; the specific heat and the macronodes are functions of $\varrho_{\tau_{\mathrm L}}$. The spectrum of $\varepsilon^\Pi I+L^\Pi$ is that of $L^\Pi$ shifted by $\varepsilon^\Pi$, so a cutoff $1/\tau_{\mathrm L}<\varepsilon^\Pi$ applied to it keeps no mode. With uniform leak $\xin(C)=\varepsilon^\Pi+\lambda_2(L_C)$ and $\xout(C)=\varepsilon^\Pi+\cut(C)/\lvert C\rvert$, so $r(C)\ge r_0$ reads $\varepsilon^\Pi(r_0-1)\le\lambda_2(L_C)-r_0\cut(C)/\lvert C\rvert$.
\end{proof}

\begin{proof}[Proof of \cref{prop:ren-lrg-agreement}]
For $i\in C_1$, $(L^\Pi u_-)_i=\sum_{j\in C_2}a^\Pi_{ij}\cdot2/N^{1/2}=2\bar a^\partial(u_-)_i$ with $u_-=(\one_{C_1}-\one_{C_2})/N^{1/2}$, and likewise on $C_2$, so $u_-$ is an eigenvector with eigenvalue $2\bar a^\partial$. The span of $\one$ and $u_-$ is invariant, hence so is its complement, the vectors orthogonal to both block indicators, on which $L^\Pi\succeq L_{C_1}\oplus L_{C_2}\succeq g_{\min}$. With unit eigenvectors $\mathsf v_k$, $(\varrho_{\tau_{\mathrm L}})_{ij}\tr e^{-\tau_{\mathrm L}L^\Pi}=N^{-1}(1\pm e^{-\tau_{\mathrm L}\lambda_2})+\sum_{k\ge3}e^{-\tau_{\mathrm L}\lambda_k}\mathsf v_k(i)\mathsf v_k(j)$, with the sign $+$ within blocks and $-$ across, and by the Cauchy--Schwarz inequality the last sum is at most $e^{-\tau_{\mathrm L}\lambda_3}$ in absolute value. Every pair within a block therefore exceeds every pair across in absolute value when $2e^{-\tau_{\mathrm L}\lambda_2}/N>2e^{-\tau_{\mathrm L}\lambda_3}$, that is $\tau_{\mathrm L}(\lambda_3-\lambda_2)>\ln N$. Exactly two modes are kept, the mode $\lambda_3$ being discarded also when $1/\tau_{\mathrm L}=\lambda_3$ because the rule of \cref{sec:ren-network} is strict, and joining pairs in descending order completes both blocks before any pair across, stopping at the two blocks. Taking $\tau_{\mathrm L}$ just below $1/(2\bar a^\partial)$, the conditions hold when $\lambda_3/(2\bar a^\partial)>1+\ln N$, for which $g_{\min}/\bar a^\partial>2(1+\ln N)$ suffices. The contrast statements are \cref{thm:col-bridge}(i) and (iv), with $\lambda_2(L^\Pi)=2\bar a^\partial$ exactly for an equitable cut.
\end{proof}

\begin{proof}[Proof of \cref{thm:ren-closure}]
(a) \Cref{lem:ren-galerkin}, with $\Vert\Pi_{\perp\Phi}\Vert^2/\varkappa_\perp\le\Vert\Pi_{\perp\Phi}\Vert^2/\xin^{\min}\le\sigma^2_{\mathrm{cut}}/\xin^{\min}$ and $\varkappa_\perp\ge\xin^{\min}$ from \cref{thm:ren-memory}(d).

(b) $\bar e=\diag(\lvert X\rvert)^{-1}\mathsf Pe$ and $e=\mathsf P^\top\bar e+e_\perp$ with $\mathsf Pe_\perp=0$, so $\bar e(t+1)=\diag(\lvert X\rvert)^{-1}\mathsf P(We(t)+\zeta(t))=W^{(1)}\bar e(t)+\diag(\lvert X\rvert)^{-1}\mathsf PWe_\perp(t)+\bar\zeta(t)$. The entries of $W^{(1)}$ are $\lvert X\rvert^{-1}\sum_{x\in X,y\in Y}W_{xy}$ and its row sums are $\lvert X\rvert^{-1}\sum_{x\in X}(W\one)_x$. Under \textup{(B)} the local noises are independent across agents with variances at most $q$, so the $\bar\zeta_X=\lvert X\rvert^{-1}\sum_{x\in X}\zeta_x$, averages over disjoint sets, are independent across blocks with variances $\lvert X\rvert^{-2}\sum_{x\in X}\Var(\zeta_x)\le q/\lvert X\rvert$, with equality when every variance equals $q$, and $\bar e$ is a linear function of $e$. The memory term vanishes for every $e_\perp\in\ker\mathsf P$ if and only if $\mathsf PW\ker\mathsf P=0$, that is $\mathsf PW=\bar W\mathsf P$ for some $\bar W$, which is lumpability; then $\bar e(t+1)=\diag(\lvert X\rvert)^{-1}\bar W\diag(\lvert X\rvert)\bar e(t)+\bar\zeta(t)$.

(c) For blocks of size $s$, $\Phi\one_S/\lvert S\rvert^{1/2}=\one_{\bigcup S}/\lvert\bigcup S\rvert^{1/2}$, so the outward rate of $S$ in $\Pi_\Phi$ is $\xout(\bigcup S)$. With $\Phi_S$ the block frame of the blocks of $S$, $(\Phi_S^\top L_{\bigcup S}\Phi_S)_{XY}=-\cut(X,Y)/s$ for $X\neq Y$ and $(\Phi_S^\top L_{\bigcup S}\Phi_S)_{XX}=\sum_{Y\in S\setminus\{X\}}\cut(X,Y)/s$, and the coarse leaks equal $\varepsilon^\Pi$, so the internal operator of $S$ is $\varepsilon^\Pi I+\Phi_S^\top L_{\bigcup S}\Phi_S$. Hence $\xin^{(1)}(S)$ is the minimum of the Rayleigh quotient of $\varepsilon^\Pi I+L_{\bigcup S}$ over the vectors of $\one_{\bigcup S}^\perp$ that are constant on blocks, a subspace of the set over which $\xin(\bigcup S)$ is the minimum; the two minima agree exactly when the subspace contains a minimizer. Contrast is invariant under scaling, so the coarse medium $s\Pi_\Phi$ gives the same values, with uniform leak. The collective statement is \cref{def:col-collective} applied to the coarse medium. Under the exchangeable clause the fine baseline is $1+\frac N{N-1}\bar d/\varepsilon^\Pi$ (\cref{prop:ren-exchangeable}); the coarse medium in the orthonormal convention has $M=N/s$ units, leak $\varepsilon^\Pi$ and mean total coupling per unit $M^{-1}\sum_X\cut(X)/s=\bar c$, so its baseline is $1+\frac M{M-1}\bar c/\varepsilon^\Pi$, and $\frac{M(N-1)}{(M-1)N}=\frac{N-1}{N-s}$. The nested statement is \cref{thm:ren-nested}(c).

(d) \Cref{prop:ren-codes,prop:ren-code-levels,prop:ren-literal-levels}.

(e) Averaging the record equation of \cref{def:fw-record-model} over $X$ gives the coarse record equation. The read term of the agents' equation there averages to $\beta_k\lvert X\rvert^{-1}\sum_{x\in X}\sum_y\mathsf K_k(x-y)z_{k,y}$, which is $\beta_k\bar z_{k,X}$ for local records and $\beta_k\sum_Y\mathsf K^{(1)}_k(X,Y)\bar z_{k,Y}$ when the record field is constant on blocks; the self term averages exactly and the message term is that of (b); record noises are independent across sites. On $\Z^D$ with cubic blocks the mean of the block averages is the mean of all agents, and likewise for the records, so the mode of wavevector zero, which obeys the reduced model of \cref{prop:sca-records}, is the same at every level.

(f) \Cref{ex:ren-no-contrast,ex:ren-block-average,lem:met-discrete-closure,prop:ren-exchangeable,prop:ren-per-level}.
\end{proof}
